\scanpage{108}
\mspage{57}
qui implique déjà leur $W$-lissité sur $S$), et $f$ compatible avec les scindages et a fortiori cartésien. Énoncé dual pour les catégories cofibrées \inserted{scindées} et les foncteurs str. cocartésiens entre elles qui sont $W$-propres fibre par fibre : ils sont $W$-propres.

2°) J'étais intéressé, dans un cas de catégories fibrées scindées \inserted{$X,Y$ sur $S$} et d'un foncteur strictement cartésien $f:X\unsure{\relbar|\longrightarrow}Y$, supposé $W$-parfait fibre par fibre, de conclure que $f$ soit aussi $W$-parfait. \deleted{\unsure{Je savais} \unsure{Comme} \unsure{par 1°)}} (C'\deleted{\unsure{était}}\unsure{est|était} un cas où de plus on sait \deleted{\unsure{déjà}} que
$X,Y$ sont aussi des $W$-fibrés sur $S$, i.e. les foncteurs chgt de base $X_{s'}\longrightarrow X_s$, $Y_{s'}\longrightarrow Y_s$ sont dans $W$. Mais ça n'a pas beaucoup servi \ldots) La question qui restait, vu \uline{1°}, était celle de la $W$-\uline{propreté} de $f$. Mais c'est sans espoir, \inserted{(même pour $Y=S$)} en vertu de \unsure{la|ce} \unsure{brillante|brillant}

\uline{Proposition 2} Supposons que
\diagram{\begin{tikzcd}[ampersand replacement=\&, column sep=small,row sep=tiny]
\& a \\
S=\unsure{o|0}\arrow[ur]\arrow[dr] \& {} \\
\& b
\end{tikzcd}}
, et soit $\deleted{\unsure{X}}Z$ une catégorie sur $S$\deleted{\unsure{.}} \deleted{(prenons $Y=S$)}. Alors $\Phi(Z)$ est propre sur $S$ sss $X=\varnothing$.

Prenons alors dans 2° ci-dessus $X=\Phi(Z)$, $Y=S$, $f:\Phi(Y)\longrightarrow S$ le morphisme structural (donc les $f_s$ sont $W$-parfaits), alors $f$ n'est propre que si $X=\varnothing$. (Pourtant, si on prend $\deleted{\unsure{X}}Z$ $W$-fibré sur $S$, $\Phi(Z)$ l'est aussi, et on a un \unsure{contreexple} à 2°).

Mais on \unsure{s'est|c'est} quand même convaincu du résultat suivant :

\uline{Prop. 3} Soient $X,Y$ cat. fibrées scindées sur $S$, $f:X\longrightarrow Y$ un foncteur str. cartésien, tel que
% archive sheet number: 108
